\documentclass[10pt]{extarticle}
\usepackage[a4paper,margin=1.5cm]{geometry}
\usepackage[T1]{fontenc}
\usepackage{amsmath,amssymb}
\usepackage{graphicx}
\usepackage[export]{adjustbox}
\usepackage{multicol}
\usepackage{enumitem}
\usepackage{booktabs,array,tabularx}
\setlength{\columnsep}{0.7cm}
\setlength{\parindent}{0pt}
\setlength{\parskip}{3pt}
\pagestyle{empty}
\makeatletter
% figure registry (figdefs.tex is generated by the extract step)
\newcommand{\qfigdef}[4]{\@namedef{qf@file@#1}{#2}\@namedef{qf@w@#1}{#3}\@namedef{qf@wide@#1}{#4}}
\newcommand{\qf@check}[1]{\@ifundefined{qf@file@#1}{\PackageError{score2}{Unknown figure id #1}{}}{}}
\newcommand{\qf@img}[4][]{\includegraphics[width=#2\dimexpr\@nameuse{qf@w@#4}pt\relax,max width=#3,max height=0.45\textheight,#1]{\@nameuse{qf@file@#4}}}
\newif\ifqf@full
\newif\ifqf@boxed
\newsavebox{\qf@box}
\newcommand{\qf@open}{\global\qf@fulltrue\par\global\qf@boxedtrue
  \setbox\qf@box\vbox\bgroup\hsize\textwidth\linewidth\textwidth}
\newcommand{\qf@emit}[1]{\par\egroup\global\qf@boxedfalse
  \ifdim\dimexpr\ht\qf@box+\dp\qf@box\relax<#1\textheight\noindent\usebox{\qf@box}\par\else\unvbox\qf@box\fi}
\newcommand{\qfullsplit}{\ifqf@boxed\qf@emit{0.85}\fi}
\newcommand{\qf@close}{\ifqf@boxed\qf@emit{0.45}\fi\global\qf@fullfalse}
\newcommand{\qfullbegin}{\end{multicols}\qf@open}
\newcommand{\qfullend}{\qf@close\begin{multicols}{2}}
\newcommand{\qfullstart}{\qf@open}
\newcommand{\qfulldone}{\qf@close}
\newcommand{\qfig}[2][1]{\qf@check{#2}%
  \ifqf@full\qf@blk{#1}{#2}\else
  \ifnum\@nameuse{qf@wide@#2}=1 \expandafter\@firstoftwo\else\expandafter\@secondoftwo\fi
  {\end{multicols}\begin{center}\qf@img{#1}{\textwidth}{#2}\end{center}\begin{multicols}{2}}%
  {\qf@blk{#1}{#2}}\fi}
\newcommand{\qf@blk}[2]{\par\smallskip{\centering\qf@img{#1}{\linewidth}{#2}\par}\smallskip}
\newcommand{\qfigi}[2][1]{\qf@check{#2}\qf@img[valign=c]{#1}{\linewidth}{#2}}
\makeatother
\newcommand{\papersec}[1]{\par\bigskip{\centering\Large\bfseries PAPER-#1\par}\medskip}
\newcommand{\partsec}[2]{\par\medskip\noindent{\large\bfseries PART-#1: #2}\par\smallskip}
\newcommand{\sect}[1]{\par\smallskip\noindent\textbf{#1}\par}
\newcommand{\qn}[1]{\par\smallskip\noindent\textbf{#1.}~}
\newcommand{\paraq}[2]{\par\smallskip\noindent\textbf{Paragraph for Questions #1 and #2}\par}
\newenvironment{qoptions}{\begin{itemize}[leftmargin=1.8em,itemsep=1pt,topsep=2pt,label={}]}{\end{itemize}}
\newcommand{\opt}[2]{\item[(#1)] #2}
% figure registry (generated by score2_tools.py finalize)
\qfigdef{p1_phy_s1_q02_1}{Score-II - Test 47 (30 Mar 2025)-img/p1_phy_s1_q02_1.png}{311.5}{1}
\qfigdef{p1_phy_s1_q06_1}{Score-II - Test 47 (30 Mar 2025)-img/p1_phy_s1_q06_1.png}{49.7}{0}
\qfigdef{p1_phy_s1_q13_1}{Score-II - Test 47 (30 Mar 2025)-img/p1_phy_s1_q13_1.png}{173.6}{0}
\qfigdef{p1_phy_s1_q14_6}{Score-II - Test 47 (30 Mar 2025)-img/p1_phy_s1_q14_6.png}{166.7}{0}
\qfigdef{p1_phy_s1_q16_1}{Score-II - Test 47 (30 Mar 2025)-img/p1_phy_s1_q16_1.png}{141.6}{0}
\qfigdef{p1_phy_s1_q18_1}{Score-II - Test 47 (30 Mar 2025)-img/p1_phy_s1_q18_1.png}{185.3}{0}
\qfigdef{p1_phy_s1_q20_1}{Score-II - Test 47 (30 Mar 2025)-img/p1_phy_s1_q20_1.png}{129.2}{0}
\qfigdef{p1_phy_s2_q01_2}{Score-II - Test 47 (30 Mar 2025)-img/p1_phy_s2_q01_2.png}{135.3}{0}
\qfigdef{p1_phy_s2_q01_3}{Score-II - Test 47 (30 Mar 2025)-img/p1_phy_s2_q01_3.png}{180.0}{0}
\qfigdef{p1_phy_s2_q04_1}{Score-II - Test 47 (30 Mar 2025)-img/p1_phy_s2_q04_1.png}{430.5}{1}
\qfigdef{p1_phy_s2_q05_1}{Score-II - Test 47 (30 Mar 2025)-img/p1_phy_s2_q05_1.png}{272.7}{1}
\qfigdef{p1_chem_s1_q09_1}{Score-II - Test 47 (30 Mar 2025)-img/p1_chem_s1_q09_1.png}{186.7}{0}
\qfigdef{p1_chem_s1_q12_1}{Score-II - Test 47 (30 Mar 2025)-img/p1_chem_s1_q12_1.png}{303.9}{1}
\qfigdef{p1_chem_s1_q12_2}{Score-II - Test 47 (30 Mar 2025)-img/p1_chem_s1_q12_2.png}{274.0}{1}
\qfigdef{p1_chem_s1_q17_1}{Score-II - Test 47 (30 Mar 2025)-img/p1_chem_s1_q17_1.png}{72.0}{0}
\qfigdef{p1_chem_s1_q17_2}{Score-II - Test 47 (30 Mar 2025)-img/p1_chem_s1_q17_2.png}{55.6}{0}
\qfigdef{p1_chem_s1_q17_3}{Score-II - Test 47 (30 Mar 2025)-img/p1_chem_s1_q17_3.png}{76.5}{0}
\qfigdef{p1_chem_s1_q19_1}{Score-II - Test 47 (30 Mar 2025)-img/p1_chem_s1_q19_1.png}{216.9}{0}
\qfigdef{p1_chem_s1_q20_2}{Score-II - Test 47 (30 Mar 2025)-img/p1_chem_s1_q20_2.png}{69.2}{0}
\qfigdef{p1_chem_s1_q20_3}{Score-II - Test 47 (30 Mar 2025)-img/p1_chem_s1_q20_3.png}{130.3}{0}
\begin{document}
\begin{center}{\Large\bfseries JEE (Main+Advanced) Score-II --- Test 47 (30 Mar 2025)}\end{center}
\begin{multicols}{2}

\partsec{1}{PHYSICS}
\sect{Section-I}

\qn{1} Which of the following statements is incorrect?
\begin{qoptions}
\opt{A}{Field obtained by the Superposition of two conservative fields is always conservative.}
\opt{B}{Field obtained by the Superposition of two non conservative fields is always non-conservative}
\opt{C}{Field obtained by the Superposition of two non-conservative fields can be non-conservative field}
\opt{D}{Field obtained by the Superposition of a conservative \& non-conservative field will always be non-conservative field}
\end{qoptions}

\qfullbegin
\qn{2} All the capacitors were uncharged before they are connected in circuit. Now the charge on the 4$\mu$F capacitor is
\qfig{p1_phy_s1_q02_1}
\qfullsplit
\begin{qoptions}
\opt{A}{Zero}
\opt{B}{$4\mu C$}
\opt{C}{$8\mu C$}
\opt{D}{$12\mu C$}
\end{qoptions}
\qfullend

\qn{3} \textbf{ASSERTION:} Electric current is a vector quantity\\
\textbf{REASON:} It has both magnitude \& direction
\begin{qoptions}
\opt{A}{Both Assertion and Reason are true and Reason is the correct explanation of Assertion.}
\opt{B}{Both Assertion and Reason are true but Reason is not the correct explanation of Assertion.}
\opt{C}{Assertion is true but Reason is false}
\opt{D}{Assertion is false but Reason is true}
\end{qoptions}

\qn{4} When two identical voltmeters are connected in series to an emf source the readings of each of voltmeter is 6 V. When both of the Voltmeters are connected in parallel to the same emf source their reading is 10 V. Find emf (in Volt).
\begin{qoptions}
\opt{A}{8}
\opt{B}{$\dfrac{90}{7}$}
\opt{C}{12}
\opt{D}{30}
\end{qoptions}

\qn{5} A body of mass m rises to a height h = R/10 from the surface of earth, where R is the radius of earth. If g is the acceleration due to gravity at the surface of earth the increase in potential energy is
\begin{qoptions}
\opt{A}{$\dfrac{1}{11}mgh$}
\opt{B}{$\dfrac{10}{11}mgh$}
\opt{C}{$\dfrac{11}{12}mgh$}
\opt{D}{$mgh$}
\end{qoptions}

\qn{6} Two blocks of masses m and nm are connected by a massless string passing over a frictionless pulley. The value of n for which both the blocks moves with an acceleration of g/6 is
\qfig{p1_phy_s1_q06_1}
\begin{qoptions}
\opt{A}{5/7}
\opt{B}{7/5}
\opt{C}{Both (A) and (B)}
\opt{D}{None}
\end{qoptions}

\qn{7} Two steel wires having same length are suspended from a ceiling under the same load. If the ratio of their energy stored per unit volume $u_1/u_2$ is 1 : 9, the ratio of their diameters $d_1/d_2$ is
\begin{qoptions}
\opt{A}{$1:\sqrt{3}$}
\opt{B}{$3:1$}
\opt{C}{$\sqrt{3}:1$}
\opt{D}{$1:3$}
\end{qoptions}

\qn{8} Force acting on a particle moving in a straight line varies with the velocity of the particle as $F=\dfrac{\alpha}{vt}$ where $\alpha$ is constant. The work done by this force from time $t_1$ to $t_2$ is :-
\begin{qoptions}
\opt{A}{$\alpha\ln\dfrac{t_2}{t_1}$}
\opt{B}{$\alpha\ln\dfrac{t_1}{t_2}$}
\opt{C}{$2\alpha(t_2-t_1)$}
\opt{D}{$\alpha^2(t_2-t_1)$}
\end{qoptions}

\qn{9} List-I gives same incomplete nuclear reactions. Match them with the reaction in List-II.

\begin{center}
\resizebox{\linewidth}{!}{%
\begin{tabular}{|c|l|c|l|}
\hline
 & \textbf{List-I} & & \textbf{List-II} \\
\hline
P & $^{237}\mathrm{Np}\longrightarrow\,^{233}\mathrm{Pa}+\_\_\_$ & 1 & $\beta^+$ decay \\
\hline
Q & $^{228}\mathrm{Ra}\longrightarrow\,\_+\_+\bar\nu$ (antineutrino) & 2 & $\beta^-$ decay \\
\hline
R & $^{41}\mathrm{Ca}\longrightarrow\,\_+\_+\nu$ (neutrino) & 3 & $\alpha$ decay \\
\hline
S & $^{197}_{\ 80}\mathrm{Hg}^*\longrightarrow\,^{197}_{\ 80}\mathrm{Hg}+\_$ & 4 & $\gamma$ decay \\
\hline
\end{tabular}}
\end{center}
\begin{qoptions}
\opt{A}{P $\to$ 4; Q $\to$ 2; R $\to$ 1; S $\to$ 3}
\opt{B}{P $\to$ 3; Q $\to$ 2; R $\to$ 1; S $\to$ 4}
\opt{C}{P $\to$ 3; Q $\to$ 1; R $\to$ 2; S $\to$ 4}
\opt{D}{P $\to$ 4; Q $\to$ 1; R $\to$ 2; S $\to$ 3}
\end{qoptions}

\qn{10} A rod of length $\ell$ and cross-sectional area A has a variable conductivity given by $K=\alpha+\beta T$, where $\alpha,\beta$ are positive constants and T is temperature (in Kelvin) of the isothermal surface. Two ends of the rod are maintained at temperatures $T_1$ and $T_2$ ($T_1>T_2$). Heat current flowing through the rod will be
\begin{qoptions}
\opt{A}{$\dfrac{A(T_2-T_1)}{\ell}\left[\beta+\dfrac{\alpha}{2}(T_2+T_1)\right]$}
\opt{B}{$\dfrac{A(T_2-T_1)}{\ell}\left[\beta-\dfrac{\alpha}{2}(T_2+T_1)\right]$}
\opt{C}{$\dfrac{A(T_2-T_1)}{\ell}\left[\alpha-\dfrac{\beta}{2}(T_2+T_1)\right]$}
\opt{D}{$\dfrac{A(T_2-T_1)}{\ell}\left[\alpha+\dfrac{\beta}{2}(T_2+T_1)\right]$}
\end{qoptions}

\qn{11} Which of the following phenomenon confirms the transverse nature of light?
\begin{qoptions}
\opt{A}{Interference}
\opt{B}{Diffraction}
\opt{C}{Polarization}
\opt{D}{Scattering}
\end{qoptions}

\qn{12} A cubical room of side 5m has a uniform magnetic field B. Total magnetic energy in the room will be.
\begin{qoptions}
\opt{A}{$\dfrac{25B^2}{2\mu_0}$}
\opt{B}{$\dfrac{125B^2}{2\mu_0}$}
\opt{C}{$\dfrac{B^2}{50\mu_0}$}
\opt{D}{$\dfrac{B^2}{250\mu_0}$}
\end{qoptions}

\qn{13} If the liquid is pushed \& released, the time period of oscillations will be (length of liquid column = $\ell$) (assume the cross sectional area of both sides of the tube to be same)
\qfig{p1_phy_s1_q13_1}
\begin{qoptions}
\opt{A}{$2\pi\sqrt{\dfrac{\ell}{g(\sin\alpha-\sin\beta)}}$}
\opt{B}{$2\pi\sqrt{\dfrac{\ell}{g(\sin\alpha+\sin\beta)}}$}
\opt{C}{$2\pi\sqrt{\dfrac{\ell}{g(\cos\alpha+\cos\beta)}}$}
\opt{D}{$2\pi\sqrt{\dfrac{\ell}{g(\cos\alpha-\cos\beta)}}$}
\end{qoptions}

\qn{14} A rod of length L has non-uniform linear mass density which varies with distance x from left end of the rod as $\rho(x)=a+b\left(\dfrac{x}{L}\right)^2$, where a and b are constants and $0\le x\le L$. The distance of the center of mass of the rod from its left end is :-
\begin{qoptions}
\opt{A}{$\dfrac{3}{2}\left(\dfrac{a+b}{2a+b}\right)L$}
\opt{B}{$\dfrac{4}{3}\left(\dfrac{a+b}{2a+3b}\right)L$}
\opt{C}{$\dfrac{3}{4}\left(\dfrac{2a+b}{3a+b}\right)L$}
\opt{D}{$\dfrac{3}{2}\left(\dfrac{2a+b}{3a+b}\right)L$}
\end{qoptions}

\qn{15} One wall of a large container has a rectangular opening. A uniform solid cylinder of density less than that of water is used to plug that rectangular opening as shown. The cylinder is free to rotate about its horizontal fixed axis. The cylinder is kept at rest till the container is filled with water above the top of cylinder as shown and then it is released.
\qfig{p1_phy_s1_q14_6}
\textbf{Statement-1:} Just after the cylinder is released it will gain angular acceleration\\
\textbf{Statement-2:} The cylinder will acquire zero steady state kinetic energy.
\begin{qoptions}
\opt{A}{Statement 1 is true, statement-2 is true and statement-2 is correct explanation for statement-1}
\opt{B}{Statement 1 is true, statement-2 is true and statement-2 is NOT the correct explanation for statement-1}
\opt{C}{Statement 1 is true, statement-2 is false}
\opt{D}{Statement 1 is false, statement-2 is true}
\end{qoptions}

\qn{16} A uniform disc of mass m and radius R is rolling without slipping up a rough inclined plane which makes an angle 30$^\circ$ with the horizontal. If the coefficient of static and kinetic friction are each equal to $\mu$ and the only force acting on the disc are gravitational and frictional force then find the
\qfig{p1_phy_s1_q16_1}
direction and magnitude of the frictional force acting on it.
\begin{qoptions}
\opt{A}{$\dfrac{mg}{6}$ down the incline}
\opt{B}{$\dfrac{mg}{6}$ up the incline}
\opt{C}{$\dfrac{mg}{3}$ down the incline}
\opt{D}{$\dfrac{mg}{3}$ up the incline}
\end{qoptions}

\qn{17} A stone projected vertically upwards reaches up to a maximum height h. Assume the potential energy of stone at the point of projection to be zero. The ratio of its Kinetic Energy to its Potential energy at a height $\dfrac{3}{5}h$, will be.
\begin{qoptions}
\opt{A}{1 : 3}
\opt{B}{3 : 1}
\opt{C}{2 : 3}
\opt{D}{3 : 2}
\end{qoptions}

\qn{18} Carbon dioxide and Carbon monoxide molecules are placed at origin \& at a large distance along x-axis respectively as shown. Carbon monoxide can be modelled as a dipole placed along x-axis \& Carbon dioxide can be modelled as quadrupole (combination of two dipoles in opposite direction) as shown in diagram. Which of the following statement regarding electrostatic force F between Carbon dioxide and Carbon monoxide is most appropriate?
\qfig{p1_phy_s1_q18_1}
\begin{qoptions}
\opt{A}{$F\propto\dfrac{1}{x^4}$}
\opt{B}{$F=0$}
\opt{C}{$F\propto\dfrac{1}{x^5}$}
\opt{D}{force F is attractive in them.}
\end{qoptions}

\qn{19} When a metallic surface is illuminated with monochromatic light of wavelength $\lambda$, the stopping potential is $5V_0$. When the same surface is illuminated with light of wavelength $3\lambda$, the stopping potential is $V_0$. Then the work function of the metallic surface is :
\begin{qoptions}
\opt{A}{$\dfrac{hc}{6\lambda}$}
\opt{B}{$\dfrac{hc}{5\lambda}$}
\opt{C}{$\dfrac{hc}{4\lambda}$}
\opt{D}{$\dfrac{2hc}{4\lambda}$}
\end{qoptions}

\qn{20} 12 identical rods made of same material are arranged in the form of a cube. The temperature of `P' and `R' are maintained at 90$^\circ$C and 30$^\circ$C respectively. Then the temperature of point `V', when steady state is reached is :-
\qfig{p1_phy_s1_q20_1}
\begin{qoptions}
\opt{A}{65$^\circ$C}
\opt{B}{60$^\circ$C}
\opt{C}{20$^\circ$C}
\opt{D}{50$^\circ$C}
\end{qoptions}

\sect{Section-II}

\qn{1} A solid uniformly charged non-conducting sphere of radius R has a spherical cavity of radius R/3 as shown. The ratio of magnitudes of electric field at center of Cavity and that at a distance R/6 from center of Cavity is
\qfig{p1_phy_s2_q01_2}

\qn{2} \qfig{p1_phy_s2_q01_3} What will be magnitude of y - coordinate of image of point A?

\qn{3} A stone is projected horizontally at a speed of 20 m/s from the top of a high tower. Simultaneously another stone is released from the same point. Find the distance in meters between stones after 1.4 seconds

\qfullbegin
\qn{4} For the following circuit and given inputs A and B, find output between $t_5$ to $t_6$.
\qfig{p1_phy_s2_q04_1}
\qn{5} An ideal diode is connected in a circuit with resistance R = 50$\Omega$ and V = 10 volt as shown in figure, maximum value of output voltage (in V) in the positive cycle will be :
\qfig{p1_phy_s2_q05_1}
\qfullend

\partsec{2}{CHEMISTRY}
\sect{Section-I}

\qn{1} 1 mole of borax reacts with how many mole of H$_2$SO$_4$?
\begin{qoptions}
\opt{A}{1}
\opt{B}{2}
\opt{C}{4}
\opt{D}{Reaction does not take place}
\end{qoptions}

\qn{2} In which of the following C-H bond has the highest `s' character?
\begin{qoptions}
\opt{A}{acetylene}
\opt{B}{ethylene}
\opt{C}{methane}
\opt{D}{Benzene}
\end{qoptions}

\qn{3} \textbf{Assertion:} The complex $[CoCl_2(en)_2]^+$ has an octahedral geometry\\
\textbf{Reason:} The complex has noble gas E.A.N
\begin{qoptions}
\opt{A}{Both Assertion and Reason are true and Reason is the correct explanation of Assertion.}
\opt{B}{Both Assertion and Reason are true but Reason is not the correct explanation of Assertion.}
\opt{C}{Assertion is true but Reason is false.}
\opt{D}{Assertion is false but Reason is true.}
\end{qoptions}

\qn{4} The average value of C--C bond order in graphite is
\begin{qoptions}
\opt{A}{1}
\opt{B}{3/2}
\opt{C}{3/4}
\opt{D}{4/3}
\end{qoptions}

\qn{5} Which of the following specie has the bond order of 0.5
\begin{qoptions}
\opt{A}{$C_2^+$}
\opt{B}{$B_2^+$}
\opt{C}{$N_2^+$}
\opt{D}{$O_2^+$}
\end{qoptions}

\qn{6} Which of the following compound is formed in borax bead test?
\begin{qoptions}
\opt{A}{Orthoborate}
\opt{B}{Metaborate}
\opt{C}{Double oxide}
\opt{D}{Tetraborate}
\end{qoptions}

\qn{7} Weight ratio of Fe : C in $Fe_2[Fe(CN)_6]$ is : [Fe = 56]
\begin{qoptions}
\opt{A}{3/7}
\opt{B}{7/5}
\opt{C}{7/3}
\opt{D}{5/7}
\end{qoptions}

\qn{8} The degree of dissociation of HI at a particular temperature is 0.8. The volume of 2M hypo solution required to react completely with the iodine present in an equilibrium mixture of reaction, which is started by 2 mole each of H$_2$ and I$_2$ in a vessel of two litre capacity, is
\begin{qoptions}
\opt{A}{0.4 litre}
\opt{B}{0.8 litre}
\opt{C}{1.6 litre}
\opt{D}{3.2 litre}
\end{qoptions}

\qn{9} Most acidic hydrogen among the following is :
\qfig{p1_chem_s1_q09_1}
\begin{qoptions}
\opt{A}{a}
\opt{B}{b}
\opt{C}{c}
\opt{D}{d}
\end{qoptions}

\qn{10} The moles of H$^+$ from H$_2$O alone in a $1l,\ \sqrt{5}\times10^{-7}\text{ M HCl}$ solution at $25^\circ$C is $(\sqrt5=2.23)$
\begin{qoptions}
\opt{A}{$10^{-7}$}
\opt{B}{$6.85\times10^{-8}$}
\opt{C}{$3.85\times10^{-8}$}
\opt{D}{$10^{-8}$}
\end{qoptions}

\qn{11} 5 moles of an ideal gas undergoes isothermal expansion from 4.0 bar to 1.0 bar, at 27$^\circ$C. The minimum amount of heat absorbed by the gas is (ln2 = 0.70)
\begin{qoptions}
\opt{A}{2250 cal}
\opt{B}{4200 cal}
\opt{C}{4500 cal}
\opt{D}{450 cal}
\end{qoptions}

\qfullbegin
\qn{12} Alkenes form alcohols by hydration reactions in which water is added in presence of dil. H$_2$SO$_4$. Addition of water follows Markownikoff's rule. For unsymmetrical alkenes, more than one product is formed. The major product is decided by stability of the intermediate carbocation which is formed by addition of H$^+$. Alcohol formation is also done by treating alkenes with other reagents. Answer the following question on the basis of above write up.
\qfig{p1_chem_s1_q12_1}
\qfullsplit
In which reactions, the product formed is correctly given?
\begin{qoptions}
\opt{A}{I, II, III}
\opt{B}{I, II}
\opt{C}{I, III, IV}
\opt{D}{III, IV}
\end{qoptions}
\qn{13} \qfig{p1_chem_s1_q12_2}
\qfullsplit Which of the following statement is NOT correct
\begin{qoptions}
\opt{A}{Compound (I) readily undergoes elimination than compound (II)}
\opt{B}{Compound (II) readily undergoes elimination than compound (I)}
\opt{C}{Elimination in (I) and (II) easily takes place when eliminating groups are anti periplanar to each other}
\opt{D}{Product (X) forms via E2 elimination}
\end{qoptions}
\qfullend

\qn{14} The values of conductance, specific conductance, equivalent conductance and molar conductance of a definite volume of an electrolyte at a particular temperature are A mho, B S cm$^{-1}$, C S cm$^{-2}$equiv$^{-1}$ and D S cm$^{-2}$mole$^{-1}$ respectively. Which of the following will decrease on dilution?
\begin{qoptions}
\opt{A}{D}
\opt{B}{B}
\opt{C}{C}
\opt{D}{Only C and D}
\end{qoptions}

\qn{15} Match the isoelectric points with amino acids:
\begin{center}
\resizebox{\linewidth}{!}{%
\begin{tabular}{|l|l||l|l|}
\hline
\multicolumn{2}{|c||}{Amino acids} & \multicolumn{2}{c|}{Isoelectric points} \\
\hline
(W) & $H_2N-CH_2COOH$ & (I) & 9.5 \\
\hline
(X) & $HOOC-CH_2CH_2CH(NH_2)COOH$ & (II) & 6.0 \\
\hline
(Y) & $H_2N(CH_2)_4CH(NH_2)COOH$ & (III) & 3.1 \\
\hline
(Z) & $H_2NCOCH_2CH(NH_2)COOH$ & (IV) & 5.6 \\
\hline
\end{tabular}}
\end{center}
\begin{qoptions}
\opt{A}{W - II, X - III, Y - I, Z - IV}
\opt{B}{W - III, X - I, Y - IV, Z - II}
\opt{C}{W - III, X - I, Y - II, Z - IV}
\opt{D}{W - II, X - I, Y - III, Z - IV}
\end{qoptions}

\qn{16} The solubility (expressed as mole fraction of gas dissolved in final solution) of any gas in any liquid increases by :
\begin{qoptions}
\opt{A}{Increase in temperature at constant pressure}
\opt{B}{Increase in amount of solvent at same temperature and pressure}
\opt{C}{Increase in pressure at constant temperature}
\opt{D}{Increase in moles of gas above solution keeping pressure and temperature constant.}
\end{qoptions}

\qn{17} If the atomic number of M is 50. The shape of gaseous $MCl_2$ is
\begin{qoptions}
\opt{A}{\qfigi{p1_chem_s1_q17_1}}
\opt{B}{\qfigi{p1_chem_s1_q17_2}}
\opt{C}{\qfigi{p1_chem_s1_q17_3}}
\opt{D}{Cl---M---Cl (linear)}
\end{qoptions}

\qn{18} When $N_2$ goes to $N_2^+$, the N--N bond distance ........ and when $O_2$ goes to $O_2^+$ the O--O bond distance....
\begin{qoptions}
\opt{A}{increases, decreases}
\opt{B}{decreases, increases}
\opt{C}{increase in both case}
\opt{D}{decrease in both case}
\end{qoptions}

\qn{19} What is the relationship between the molecules in the following pairs?
\qfig{p1_chem_s1_q19_1}
\begin{qoptions}
\opt{A}{enantiomers}
\opt{B}{diastereomers}
\opt{C}{identical}
\opt{D}{structural isomers}
\end{qoptions}

\qn{20} Match stereo isomers for above compounds:
\begin{center}
\resizebox{\linewidth}{!}{%
\begin{tabular}{|c|c|c|c|}
\hline
\multicolumn{2}{|c|}{\textbf{Column-I}} & \multicolumn{2}{c|}{\textbf{Column-II}} \\
\hline
(A) & $CH_3\text-CH{=}CH\text-CH{=}N\text-OH$ & (P) & 2 \\
\hline
(B) & \qfigi{p1_chem_s1_q20_2} & (Q) & 4 \\
\hline
(C) & \qfigi{p1_chem_s1_q20_3} & (R) & 3 \\
\hline
(D) & \begin{tabular}{@{}c@{}}$CH_3\text-CH{=}CH\text-CH{=}CH\text-CH\text-Ph$\\ $|$\\ $CH_3$\end{tabular} & (S) & 8 \\
\hline
\end{tabular}}
\end{center}
\begin{qoptions}
\opt{A}{A $\to$ Q, B $\to$ P, C $\to$ R, D $\to$ S}
\opt{B}{A $\to$ S, B $\to$ P, C $\to$ R, D $\to$ Q}
\opt{C}{A $\to$ Q, B $\to$ R, C $\to$ P, D $\to$ S}
\opt{D}{A $\to$ P, B $\to$ P, C $\to$ Q, D $\to$ S}
\end{qoptions}

\sect{Section-II}

\qn{1} The total number of diprotic acids among the following is\\
$H_3PO_4$, $H_2SO_4$, $H_3PO_3$, $H_2CO_3$, $H_2S_2O_7$, $H_3BO_3$, $H_3PO_2$, $H_2CrO_4$, $H_2SO_3$

\qn{2} For a second order reaction, containing single reactant : A $\to$ B, it is given that :\\
$T_{75\%}=X\cdot t_{50\%}$\\
Find the value of X

\qn{3} The molecular mass of an aliphatic ketone is 86. It forms oxime with $NH_2OH$. The oxime is reduced to a resolved amine. How many structural isomer(s) of the ketone form(s) resolved amine?

\qn{4} $CH_2=C=C=C=CH_2$\\
How many maximum number atoms of the above compounds are present in one plane?

\qn{5} Total number of stereoisomerism present in complex compound $[Cr(gly)(NH_3)(H_2O)Cl_2Br]$.

\partsec{3}{MATHEMATICS}
\sect{Section-I}

\qn{1} Let $f(x)=\displaystyle\int \dfrac{x^2\,dx}{(1+x^2)\left(1+\sqrt{1+x^2}\right)}$ and $f(0)=0$. Then $f(1)$ is equal to
\begin{qoptions}
\opt{A}{$\ln\left(1+\sqrt2\right)$}
\opt{B}{$\ln\left(1+\sqrt2\right)-\dfrac{\pi}{4}$}
\opt{C}{$\ln\left(1+\sqrt2\right)+\dfrac{\pi}{4}$}
\opt{D}{None of these}
\end{qoptions}

\qn{2} A is $3\times3$ matrix whose entries belong to the set $\{-1,0,1\}$. Then the probability that A is neither symmetric nor skew-symmetric is
\begin{qoptions}
\opt{A}{$\dfrac{3^9-3^6-3^3+1}{3^9}$}
\opt{B}{$\dfrac{3^9-3^6+3^3+1}{3^9}$}
\opt{C}{$\dfrac{3^9-3^6-3^3}{3^9}$}
\opt{D}{$\dfrac{3^9-3^6+3^3}{3^9}$}
\end{qoptions}

\qn{3} Consider the function $f(x)=\sin^{-1}\left(\dfrac{2x}{1+x^2}\right)+2\tan^{-1}x$\\
\textbf{Assertion (A):} $f(10)=\pi$\\
\textbf{Reason (R):} $f(x)=\pi\ \forall\ x\ge1$
\begin{qoptions}
\opt{A}{A is true, R is true and R is correct explanation for A}
\opt{B}{A is true, R is true and R is NOT the correct explanation for A}
\opt{C}{A is true, R is false}
\opt{D}{A is false, R is true}
\end{qoptions}

\qn{4} If z is a complex number which satisfies the equation $|z-5-7i|=9$, then the maximum value of $|z-2-3i|$ is equal to
\begin{qoptions}
\opt{A}{4}
\opt{B}{9}
\opt{C}{12}
\opt{D}{14}
\end{qoptions}

\qn{5} A pair of mutually perpendicular lines is drawn from origin to form an isosceles triangle with the line $2x+3y=6$. The area of triangle thus formed is equal to
\begin{qoptions}
\opt{A}{$\dfrac{36}{13}$}
\opt{B}{$\dfrac{12}{17}$}
\opt{C}{$\dfrac{13}{5}$}
\opt{D}{$\dfrac{6}{13}$}
\end{qoptions}

\qn{6} The curve passing through (0, 1) which satisfies the differential equation $\dfrac{xdy-ydx}{xdy+ydx}=y^2\sin(xy)$, is given by
\begin{qoptions}
\opt{A}{$x+y(1+\cos(xy))=0$}
\opt{B}{$x-y(1-\cos(xy))=0$}
\opt{C}{$x+y(1-\cos(xy))=0$}
\opt{D}{$x-y(1+\cos(xy))=0$}
\end{qoptions}

\qn{7} A latus rectum subtends a right angle at the other focus of the hyperbola. The eccentricity of the hyperbola must be
\begin{qoptions}
\opt{A}{$1-\sqrt2$}
\opt{B}{$\sqrt2$}
\opt{C}{$\sqrt2-1$}
\opt{D}{$1+\sqrt2$}
\end{qoptions}

\qn{8} The first term of an Arithmetic Progression (AP) with real terms is $-1$ and the $8^{th}$ term of a Geometric Progression (GP) with real terms is 80. The $4^{th}$ terms of both progressions are equal. If the $6^{th}$ term of the GP is equal to the sum of the $6^{th}$ and $7^{th}$ terms of the AP, then the smallest possible integral value of the $9^{th}$ term of the GP is equal to
\begin{qoptions}
\opt{A}{160}
\opt{B}{2}
\opt{C}{$-160$}
\opt{D}{$-2$}
\end{qoptions}

\qn{9} Consider the following lists
\begin{center}
\resizebox{\linewidth}{!}{%
\begin{tabular}{|c|l|c|l|}
\hline
& \textbf{List-I} & & \textbf{List-II} \\
\hline
(A) & $\left\{x\in\left[-\dfrac{2\pi}{3},\dfrac{2\pi}{3}\right]:\sin x+\cos x=1\right\}$ & (I) & Has elements 2 \\
\hline
(B) & $\left\{x\in\left[-\dfrac{5\pi}{18},\dfrac{5\pi}{18}\right]:\sqrt3\tan3x=1\right\}$ & (II) & Has elements 3 \\
\hline
(C) & $\left\{x\in\left[-\dfrac{6\pi}{5},\dfrac{6\pi}{5}\right]:2\cos(2x)=\sqrt3\right\}$ & (III) & Has elements 4 \\
\hline
(D) & $\left\{x\in\left[-\dfrac{7\pi}{4},\dfrac{7\pi}{4}\right]:\sin x-\cos x=1\right\}$ & (IV) & Has elements 6 \\
\hline
\end{tabular}}
\end{center}
\begin{qoptions}
\opt{A}{A $\to$ I; B $\to$ IV; C $\to$ I; D $\to$ IV}
\opt{B}{A $\to$ I; B $\to$ I; C $\to$ IV; D $\to$ III}
\opt{C}{A $\to$ II; B $\to$ I; C $\to$ IV; D $\to$ IV}
\opt{D}{A $\to$ II; B $\to$ IV; C $\to$ I; D $\to$ III}
\end{qoptions}

\qn{10} The value of $\displaystyle\int_0^\infty\left[\dfrac{3}{1+x^2}\right]dx$ ($[.]$ denotes greatest integer function) is
\begin{qoptions}
\opt{A}{3}
\opt{B}{$\dfrac{3}{\sqrt2}$}
\opt{C}{$3\sqrt2$}
\opt{D}{infinite}
\end{qoptions}

\qn{11} The area of the region A represented by $A=\{(x,y)\mid(x,y)\in R\times R,\ 1\le|x-2|+|y+1|\le2\}$ is
\begin{qoptions}
\opt{A}{9 sq. units}
\opt{B}{4 sq. units}
\opt{C}{5 sq. units}
\opt{D}{6 sq. units}
\end{qoptions}

\qn{12} If the line $y=2-x$ is tangent to the circle S at a point P (1,1) and circle S is orthogonal to the circle $x^2+y^2+2x+2y-2=0$, then the equation of circle S is
\begin{qoptions}
\opt{A}{$x^2+y^2=2$}
\opt{B}{$x^2+y^2-3x-3y+4=0$}
\opt{C}{$x^2+y^2+x-3y=0$}
\opt{D}{$x^2+y^2-x-y=0$}
\end{qoptions}

\qn{13} A fair dice is thrown three times and the respective outcomes are named a, b and c. The probability that $f(x)=x^3+ax^2+bx+c$ is a increasing function is
\begin{qoptions}
\opt{A}{$\dfrac{5}{36}$}
\opt{B}{$\dfrac{4}{9}$}
\opt{C}{$\dfrac{1}{3}$}
\opt{D}{None of these}
\end{qoptions}

\qn{14} A focal chord of the parabola $y^2=16x$ is a tangent to the circle $(x-6)^2+y^2=2$. The length of this focal chord is equal to
\begin{qoptions}
\opt{A}{32}
\opt{B}{16}
\opt{C}{512}
\opt{D}{64}
\end{qoptions}

\qn{15} Consider the following statements about the function $f:A\to A$\\
\textbf{Statement-1:} If f is onto, then it must be one-one.\\
\textbf{Statement-2:} If f is one-one, then it must be onto.
\begin{qoptions}
\opt{A}{Only (S1) is true}
\opt{B}{Only (S2) is true}
\opt{C}{Both (S1) and (S2) is true}
\opt{D}{Both (S1) and (S2) is false}
\end{qoptions}

\qn{16} An ellipse having foci (3,1) and (1,1) and passing through the point (1,3) has eccentricity equal to
\begin{qoptions}
\opt{A}{$\sqrt2-1$}
\opt{B}{$\dfrac{\sqrt2-1}{2}$}
\opt{C}{$\sqrt3-1$}
\opt{D}{$\dfrac{\sqrt3-1}{2}$}
\end{qoptions}

\qn{17} Let the observations $x_i\,(1\le i\le10)$ satisfy the equations $\displaystyle\sum_{i=1}^{10}(x_i-5)=10$ and $\displaystyle\sum_{i=1}^{10}(x_i-5)^2=40$. If $\mu$ and $\lambda$ are the mean and the variance of the observations, $x_1-3,x_2-3,\dots,x_{10}-3$, then the ordered pair $(\mu,\lambda)$ is equal to :
\begin{qoptions}
\opt{A}{(6, 6)}
\opt{B}{(3, 6)}
\opt{C}{(6, 3)}
\opt{D}{(3, 3)}
\end{qoptions}

\qn{18} The value of $\tan70^\circ-\tan50^\circ+\tan10^\circ$ equals
\begin{qoptions}
\opt{A}{0}
\opt{B}{2}
\opt{C}{3}
\opt{D}{$\sqrt3$}
\end{qoptions}

\qn{19} Suppose $p,q,r\ne0$ and the system of equations
\[(p+a)x+by+cz=0\]
\[ax+(q+b)y+cz=0\]
\[ax+by+(r+c)z=0\]
has non-trivial solutions, then the value of $\dfrac{a}{p}+\dfrac{b}{q}+\dfrac{c}{r}$ is
\begin{qoptions}
\opt{A}{$-1$}
\opt{B}{0}
\opt{C}{1}
\opt{D}{2}
\end{qoptions}

\qn{20} If p, q, r and s are non-zero numbers such that r, s are roots of $x^2+px+q=0$ and p, q are the roots of $x^2+rx+s=0$, then the value of $(4p+q+2r+s)$ is
\begin{qoptions}
\opt{A}{0}
\opt{B}{1}
\opt{C}{2}
\opt{D}{4}
\end{qoptions}

\sect{Section-II}

\qn{1} The value of $\displaystyle\lim_{x\to2}\dfrac{2^x+2^{3-x}-6}{2^{-\frac{x}{2}}-2^{1-x}}$ is

\qn{2} If $\vec a$ and $\vec b$ are any two unit vectors, then the maximum possible integral value of $3|\vec a-\vec b|+\dfrac{7}{2}|\vec a+\vec b|$ is equal to

\qn{3} The shortest distance between the lines $\dfrac{x-1}{2}=\dfrac{y}{3}=\dfrac{z}{2}$ and $\dfrac{x-3}{2}=\dfrac{y}{5}=\dfrac{z-2}{4}$ is equal to

\qn{4} The remainder when $5^{99}$ is divided by 13 is

\qn{5} How many words with or without meaning formed using all the letters of the word ALLEN begin and ends with a consonant?

\end{multicols}
\end{document}
